\section{Results}\label{sec:results}

\subsection{Ideal-source baseline}

We first establish the ideal-pump baseline that all imperfection budgets
refer to. Figure~\ref{fig:design}(a) shows the heralded purity of
Eq.~(\ref{eq:purity}) as a function of the pump bandwidth for the 10~mm
reference crystal. The curve has the broad interior optimum expected
when the pump-envelope and phase-matching correlations of
Eq.~(\ref{eq:jsa}) are balanced \cite{grice1997,law2000,mosley2008}:
the grid-limited peak is $P_{\max} = 0.8106$ at $0.178$~THz on the
sweep's grid ($n = 512$, 6~THz frequency half-window), the top is
flat, and the canonical reference point $\sigma = 0.164$~THz gives
$P = 0.8100$ on the reference grid ($n = 256$, 4~THz half-window) used
for all Monte-Carlo studies below. The 1\% tolerance
window spans 0.159--0.202~THz, so the bandwidth itself is a forgiving
design parameter. The fall-off is nearly log-symmetric about the
optimum: equal factors from the peak give nearly equal purities
(e.g.\ $P = 0.663$ and $0.665$ a factor of $1.9$ below and above it),
and the endpoint values---$P = 0.41$ at $0.05$~THz against $0.24$ at
$1.2$~THz---differ only because the sweep extends further on the
broadband side (factors of $3.6$ and $6.7$ from the peak,
respectively). The residual
gap to $P = 1$ is the familiar penalty of an unapodized sinc
phase-matching function \cite{branczyk2011,graffitti2018}.
Figure~\ref{fig:design}(b) shows the complementary sweep: at fixed
$\sigma = 0.164$~THz the crystal length has its optimum at 10--12~mm
($P = 0.806$ at 10~mm, $0.805$ at 12~mm), falling to $0.286$ at 2~mm
and $0.51$ at 30~mm. This is the same matching condition seen from the
crystal side. The small offset between the $0.806$ quoted here and the
$0.8100$ reference value above is a frequency-window truncation effect,
not discretization: the slowly decaying sinc side lobes make the
computed purity drift with the window half-width ($P = 0.8100$,
$0.8074$, and $0.8061$ at the same operating point for the 4, 6, and
8~THz half-windows of the reference grid, the bandwidth sweep, and the
length sweep, respectively), so ideal purities are comparable only
within a fixed window and are window-independent only at the level
$P \approx 0.81$. As a narrowband anchor, a 5~GHz pump yields
$P \approx 0.045$ and $K \approx 22$ (Eq.~(\ref{eq:schmidt})) once the
grid resolves the few-GHz pump ridge (which is narrower than the pixel
of the sweep grids), recovering the strongly anti-correlated quasi-CW
limit.

\begin{figure*}[t]
\includegraphics[width=\linewidth]{fig2_design}
\caption{Ideal-source design space. (a) Heralded purity versus pump
bandwidth (log axis) for the 10~mm crystal; the dotted line marks the
grid-limited peak $P_{\max} = 0.8106$ at $0.178$~THz on a flat-topped
optimum. (b) Purity versus crystal length (2--30~mm) at the fixed
reference pump bandwidth of $0.164$~THz, with the optimum at
10--12~mm.}
\label{fig:design}
\end{figure*}

\subsection{Spectral drift and jitter}

Figure~\ref{fig:jitter} shows the heralded purity of the Monte-Carlo
mixed state (Eq.~(\ref{eq:rho})) as the pump center frequency jitters
with increasing rms amplitude. At zero jitter the ensemble reproduces
the ideal value exactly, $P(0) = 0.8100$; the seed-averaged degradation
is monotonic and reaches $P = 0.299$ at $0.5$~THz rms (individual seeds
scatter non-monotonically about this mean, with the worst seed reaching
$0.287$). The shaded band is the spread
of three independent seeds ($N = 120$ realizations each); at
$0.25$~THz rms the three seeds give $P = 0.484/0.460/0.513$, i.e.\ a
spread of about $\pm 0.03$ that dominates the negligible discretization
error ($\sim 2\times 10^{-5}$ under grid refinement at fixed window).
The seed-averaged budgets for the 10~mm
crystal are 26/61/88~GHz rms for a 1/5/10\% relative purity drop.
Re-matching the source with a 5~mm crystal and a $0.328$~THz pump
(ideal $P = 0.8183$) roughly doubles every budget to 51/121/188~GHz
(factors of 1.98--2.12; the 5~mm sweep uses a single seed): the
phase-matching bandwidth scales as $1/L$
(Eq.~(\ref{eq:pm})), so a shorter crystal simply presents a wider
acceptance to a wandering pump, at the cost of a lower pair rate that
our model does not capture.

\begin{figure}[t]
\includegraphics[width=\linewidth]{fig3_jitter}
\caption{Heralded purity versus pump center jitter/drift (rms, treated
as a quasi-static ensemble) for the 10~mm reference source (circles;
shaded band: spread of three Monte-Carlo seeds) and for a 5~mm crystal
with the pump re-matched to $0.328$~THz (squares). Dotted lines mark
the respective ideal purities.}
\label{fig:jitter}
\end{figure}

\subsection{Longitudinal-mode structure}

Multimode pump operation is far less forgiving.
Figure~\ref{fig:multimode}(a) shows the purity ladder for an
equal-weight comb at $0.25$~THz spacing:
$P = 0.810$, $0.670$, $0.535$, $0.451$, $0.370$, $0.321$, $0.297$ for
1--7 lines. A single extra mode at this spacing already costs 17\% of
the purity---more than the entire 10\% budget---so the maximum mode
count within any budget is one. Figure~\ref{fig:multimode}(b) sweeps
the spacing of a fixed 3-mode comb: at $0.02$~THz the lines are
unresolved by the JSA and $P = 0.807$, essentially the single-mode
value, while at $0.6$~THz the purity has fallen to $0.31$. The spacing
budgets are 31/76/111~GHz for 1/5/10\%. The practical reading is that a
nominally single-frequency pump that occasionally hops or lases on a
side mode is a first-order threat to heralded purity, on par with the
strongest jitter considered above.

\begin{figure*}[t]
\includegraphics[width=\linewidth]{fig4_multimode}
\caption{Longitudinal-mode structure of the pump. (a) Heralded purity
versus the number of comb lines at $0.25$~THz spacing; the dotted line
is the single-mode value. (b) Purity versus the spacing of a 3-mode
comb (0.02--0.6~THz); for spacings unresolved by the JSA the
single-mode purity is recovered.}
\label{fig:multimode}
\end{figure*}

\subsection{Intensity noise attacks a different metric}

Relative intensity noise behaves qualitatively differently
(Fig.~\ref{fig:rin}). Because a pulse-energy fluctuation rescales the
ensemble weights $w_k$ in Eq.~(\ref{eq:rho}) without changing any
normalized spectral amplitude, the heralded purity is exactly
invariant: on the $n = 64$ grid used for this study it is flat at
$0.809956$ for every RIN value from 0 to 30\% rms
(Fig.~\ref{fig:rin}(a)). What RIN does damage is the pair statistics
relevant for heralding fidelity and accidental double pairs
\cite{takeoka2015,meyerscott2020}: the double-pair boost
$\langle w^2 \rangle / \langle w \rangle^2$ follows the small-$r$
prediction $1 + r^2$, with a measured $1.0221$ at 15\% RIN against a
theoretical $1.0225$, and $1.0865$ at 30\% against $1.09$
(Fig.~\ref{fig:rin}(b)). The residual scatter about $1 + r^2$ is
finite-sample Monte-Carlo error: the boost estimator has a standard
error of $\approx 0.004$ at $r = 0.3$ for $N = 10^3$ realizations, and
the measured points fall on both sides of the theory (e.g.\ $1.0560$
against $1.0506$ at 22.5\% RIN), whereas the clipping of the weights
$w = \max[0,\, 1 + \mathcal{N}(0, r^2)]$ at zero is negligible here,
shifting the boost by only $-8\times 10^{-5}$ at $r = 0.30$. Budgeted
on the boost, a $+1\%$ increase
is reached at 10.0\% RIN and $+5\%$ at 21.7\%; $+10\%$ is not reached
within the 30\% sweep. A datasheet must therefore be read with this
split in mind: jitter and mode structure are spectral imperfections
that degrade purity, while RIN is a statistical imperfection that
leaves purity untouched and instead degrades the pair-number
statistics.

\begin{figure*}[t]
\includegraphics[width=\linewidth]{fig5_rin}
\caption{Relative intensity noise. (a) Heralded purity versus RIN
(0--30\% rms): exactly invariant, flat at $0.809956$ on the $n = 64$
grid. (b) Double-pair boost $\langle w^2 \rangle / \langle w \rangle^2$
from the Monte-Carlo ensemble ($N = 10^3$) against the small-$r$ theory
$1 + r^2$ (dashed); the dotted line marks the noiseless value of 1.}
\label{fig:rin}
\end{figure*}

\subsection{The imperfection budget}

Table~\ref{tab:budget} collects the results of
Figs.~\ref{fig:design}--\ref{fig:rin} into a single auto-generated
budget. The convention is as follows: each entry is the largest value
of the imperfection for which the heralded purity stays within a
1/5/10\% \emph{relative} drop of the ideal value of the same source
(0.8100 for the 10~mm source, 0.8183 for the re-matched 5~mm source);
the pump-bandwidth row is instead the window of bandwidths within the
stated drop of the peak, and RIN---which leaves purity exactly
invariant---is budgeted on the double-pair boost instead. The
hierarchy is unambiguous: longitudinal-mode structure is the least
tolerable imperfection, center jitter comes next with budgets in the
tens of GHz, and RIN only matters at the ten-percent level and through
a different observable.

\begin{table*}[t]
\centering
% Auto-generated by paper/make_budget.py -- do not edit by hand.
% Cells: mean +- SEM over per-seed threshold crossings (see caption).
\begin{tabular}{llrrr}
\hline\hline
Imperfection & unit & \multicolumn{3}{c}{budget for relative purity drop} \\
 & & 1\% & 5\% & 10\% \\
\hline
Center jitter rms (20\,mm)  & GHz & $13.5 \pm 0.2$ & $31.3 \pm 0.4$ & $45.3 \pm 0.6$ \\
Center jitter rms (10\,mm)  & GHz & $27.7 \pm 0.3$ & $62.7 \pm 0.7$ & $91 \pm 1$ \\
Center jitter rms (5\,mm)   & GHz & $54.1 \pm 0.9$ & $125 \pm 1$ & $181 \pm 3$ \\
Center jitter rms (10\,mm, Gaussian PM) & GHz & $30 \pm 0.3$ & $68.5 \pm 0.8$ & $102 \pm 1$ \\
Mode spacing (3-mode comb)  & GHz & $29.6 \pm 0.1$ & $75.9 \pm 0.5$ & $112.1 \pm 0.7$ \\
Additional side modes (0.25\,THz) & count & 0 & 0 & 0 \\
Pump bandwidth window       & THz & 0.159--0.202 & 0.131--0.244 & 0.114--0.281 \\
RIN (budget on pair boost)  & \%\,rms & $9.93 \pm 0.08$ & $22.3 \pm 0.2$ & beyond sweep \\
\hline\hline
\end{tabular}

\caption{Imperfection budget, auto-generated from the sweeps of
Sec.~\ref{sec:results}. Each entry is the largest imperfection keeping
the heralded purity within the stated relative drop of the ideal value
of the same source; the bandwidth row gives the tolerance window
around the purity peak, and RIN is budgeted on the double-pair boost
$\langle w^2 \rangle / \langle w \rangle^2$ because the purity is
exactly RIN-invariant.}
\label{tab:budget}
\end{table*}
